Satellite embeddings have become a practical interface between large-scale Earth observation data and downstream geospatial analysis, yet the literature is still organized mainly around foundation models rather than the embeddings they produce. This review reframes the area from an embedding-centered perspective. We first define satellite embeddings as reusable latent representations derived from satellite or multimodal Earth observation data, and organize existing methods along three axes: model family, data and modality support, and deployment mode. We then consolidate \modelcount{} representative models into a unified dossier covering publication source, authorship, spatial scale, temporal semantics, architecture, pretraining signal, embedding dimensionality, runtime behavior, and code or artifact availability. Next, we outline an application-oriented benchmark agenda spanning classification, segmentation, regression, and retrieval, with a small-area comparative protocol that can be executed across both on-the-fly encoders and precomputed embedding products. Finally, we describe an interactive system layer that exposes these models through a lightweight explorer for region selection, model comparison, and rapid downstream experimentation. The resulting review is intended to serve simultaneously as a technical reference, an empirical evaluation framework, and a reproducible software entry point for the satellite embedding community.
